% ============================================================================
%   CHAPTER 1 — INTRODUCTION
%   Aivancity: 3-5 pages, funnel structure (general → specific)
% ============================================================================

\chapter{Introduction}
\label{chap:introduction}

% ----------------------------------------------------------------------------
\section{Macro Context: AI in Healthcare}
\label{sec:intro-macro}

% [TODO — 0.5-1 page]
% - AI transformation of healthcare
% - The specific role of LLMs since 2022
% - Societal stakes: efficiency, accessibility, safety, trust

% ----------------------------------------------------------------------------
\section{Micro Context: Clinical Trial Matching}
\label{sec:intro-micro}

% [TODO — 0.5-1 page]
% - The clinical trial matching problem: what it is, why it matters
% - Current bottleneck: manual review, low enrollment rates
% - LLM-based approaches: TrialGPT (Jin et al. 2024) as the reference

% ----------------------------------------------------------------------------
\section{Problem Statement}
\label{sec:intro-problem}

% [TODO — 0.5 page]
% - LLMs produce eligibility scores, but ARE THESE SCORES RELIABLE?
% - The specific problem of calibration
% - Why calibration matters clinically (Van Calster 2019, TRIPOD+AI)
% - The gap: no systematic audit of LLM calibration in trial matching

% ----------------------------------------------------------------------------
\section{Research Question and Hypotheses}
\label{sec:intro-question}

% [TODO — 0.5-1 page]
% Main research question:
%   "Are the eligibility scores produced by LLM-based patient-trial matching
%    systems probabilistically reliable, and does RLHF post-training affect
%    this reliability?"
%
% Pre-registered falsifiable hypotheses:
%   H1: Current LLM pipelines are poorly calibrated (ECE > 0.05).
%   H2: RLHF post-training degrades calibration
%       [ECE(Instruct) > ECE(Base), permutation test with Bonferroni].
%   H3: LLM scores exhibit ordinal validity (C-index > 0.70).
%   H4: Selective abstention improves clinical quality
%       [F1@80% coverage > F1@100% + 0.02].

% ----------------------------------------------------------------------------
\section{Contributions}
\label{sec:intro-contributions}

% [TODO — 0.5 page]
% List the 9 original contributions:
%   1. Systematic calibration audit of 12 LLM configurations, with
%      cross-validation and comparison against naïve baselines
%   2. Empirical demonstration that RLHF significantly degrades calibration
%      across two independent model families (H2, main pre-registered
%      contribution, extending Kadavath 2022 to the clinical domain)
%   3. The enriched-prompt paradox: expert prompts degrade calibration in
%      5 out of 6 models
%   4. Systematic inclusion/exclusion asymmetry in LLM predictions
%   5. Model-specific evasion biases: GPT-4 over-uses "not applicable",
%      open-source LLMs over-use "not enough information"
%   6. Post-hoc isotonic recalibration brings 11 of 12 configurations below
%      the calibration threshold, robustly across cross-validation folds
%   7. Discovery of mode collapse: 5 of 12 configurations predict "not
%      enough information" in more than 97% of cases
%   8. Structural nature of mode collapse: forcing decisive labels does not
%      restore reasoning
%   9. Complete failure of medical LLMs on truly discriminative cases,
%      hidden by aggregate metrics

% ----------------------------------------------------------------------------
\section{Structure of the Thesis}
\label{sec:intro-structure}

% [TODO — 0.5 page]
% One sentence per chapter:
%   Chapter 2 reviews the state of the art on clinical trial matching,
%     LLM calibration, and the effect of RLHF on model behavior.
%   Chapter 3 describes the project context, the dataset used, and the
%     scientific constraints.
%   Chapter 4 details the experimental methodology, including model
%     selection, prompt design, metrics, and statistical analysis plan.
%   Chapter 5 presents the results structured around the four
%     pre-registered hypotheses and five original findings.
%   Chapter 6 discusses the implications, limitations, and future work,
%     including a proposed multi-agent debate extension.
